\section{Introduction}

Under EU REACH\cite{EU_REACH_2006} and K-REACH,\cite{K_REACH_2018} industrial chemical 
manufacturers and importers must assess genotoxic potential as part of their registration 
dossiers; in silico QSAR predictions are explicitly accepted as supporting evidence when 
generated by validated models with sufficient applicability domain 
coverage.\cite{ECHA_QSAR_2008}

For pharmaceuticals, the ICH M7(R2) guideline plays an analogous role by endorsing 
statistical QSAR alongside rule-based expert systems for assessing the mutagenic potential 
of impurities.\cite{ICH_M7_2017} Across both regulatory frameworks, the Ames bacterial 
reverse mutation assay\cite{Ames_1973,Mortelmans_2000} is the central predictive target, 
with in-vitro chromosome aberration (CA) and in-vivo micronucleus 
(MN)\cite{OECD_473,OECD_474} serving as higher-tier endpoints. Reliable models could 
reduce animal testing in line with the 3Rs principle.\cite{Russell_Burch_1959}

Published Ames QSAR models report a wide spread of metrics---accuracies of 0.80--0.90 and 
MCC of 0.60--0.80\cite{Hansen_2009,Xu_2012,Honma_2019,Li_2021}---but these come from 
different splits, datasets, preprocessing pipelines, and protocols, making direct comparison 
difficult. At least three factors can inflate apparent performance.

\textbf{Splitting strategy.} Random splits allow structurally similar compounds (sharing the 
same Murcko scaffold\cite{Bemis_Murcko_1996}) to appear in both training and test set, 
producing optimistic estimates. Scaffold-based 
splitting\cite{Sheridan_2013,Wu_MoleculeNet_2018} is more realistic, but the magnitude of 
the random-vs-scaffold gap has not been systematically quantified across genotoxicity 
endpoints.

\textbf{Source mixing.} Most large Ames datasets aggregate compounds from pharmaceutical 
libraries and industrial chemical registries,\cite{Hansen_2009,Honma_2019} which differ in 
composition, prevalence, and likely in data quality and labeling conventions. When mixed, 
models may learn source membership rather than toxicity mechanisms---a form of data source 
heterogeneity noted\cite{Tropsha_2010,Hanser_2016} but never systematically measured. One 
component of this heterogeneity is prior shift: the large prevalence gap between sources can 
miscalibrate the default decision threshold independently of any structural difference. 
However, threshold decomposition analysis presented here shows that prior shift accounts for 
only part of the observed performance gap, with a residual that no recalibration can 
eliminate.

\textbf{Preprocessing.} Salt stripping, metal removal, and charge normalization can 
substantially alter molecular representations,\cite{Fourches_2010,Fourches_2016} yet these 
steps are often undocumented or applied uniformly across endpoints.

Most studies evaluate a single endpoint with a single split, focus on algorithm 
comparison,\cite{Yang_2019,Mayr_2016} rarely test cross-source generalization, and seldom 
report statistical significance.\cite{Demsar_2006} The relative contributions of split 
method, source composition, and preprocessing have not been disentangled in a single 
consistent framework.

This study addresses these gaps by:
\begin{enumerate}
    \item Introducing an evaluation cascade---scaffold-split test, random vs.\ scaffold CV, 
    and cross-source leave-domain-out---applied consistently to a single locked 
    configuration.

    \item Showing that, on the Ames benchmark, data source heterogeneity drives a 
    scaffold-CV-to-LDO MCC drop of 0.390---an order of magnitude larger than the split-method 
    effect ($\Delta = 0.046$)---and that threshold-tuning analysis attributes only part of 
    this drop to prior shift, leaving a residual that no recalibration can eliminate.

    \item Finding that in-domain algorithm differences are small (MCC range = 0.063) but 
    that cross-source robustness varies sharply (RF 0.325 vs.\ XGB 0.122 in the harder 
    direction), so algorithm choice matters substantially more under source shift than 
    in-domain comparisons suggest.

    \item Reporting endpoint-specific preprocessing effects, with salt stripping improving 
    the best in-vivo model by $+0.22$ MCC while harming in-vitro.

    \item Releasing an open-source pipeline for full reproducibility.
\end{enumerate}
